\documentclass[11pt]{article}
\usepackage[margin=1in]{geometry}
\usepackage[T1]{fontenc}
\usepackage[utf8]{inputenc}
\usepackage{lmodern}
\usepackage{microtype}
\usepackage[table]{xcolor}
\usepackage{amsmath,amssymb,mathtools,bm}
\usepackage{booktabs,array,longtable,tabularx}
\usepackage{hyperref}
\usepackage{enumitem}
\usepackage{fancyhdr}
\usepackage{titlesec}
\usepackage{tcolorbox}
\usepackage{ragged2e}
\usepackage{setspace}
\hypersetup{colorlinks=true,linkcolor=blue!45!black,urlcolor=blue!45!black,citecolor=blue!45!black}
\definecolor{TitleBlue}{HTML}{163A5F}
\definecolor{AccentTeal}{HTML}{2D6F73}
\definecolor{SoftBlue}{HTML}{EAF3F8}
\definecolor{SoftTeal}{HTML}{EDF8F6}
\definecolor{SoftGray}{HTML}{F5F7FA}
\pagestyle{fancy}
\fancyhf{}
\lhead{PINN\_M3DC1}
\rhead{Project Plan}
\cfoot{\thepage}
\fancypagestyle{plain}{%
  \fancyhf{}%
  \lhead{PINN\_M3DC1}%
  \rhead{Project Plan}%
  \cfoot{\thepage}%
}
\setlength{\headheight}{14pt}
\setlength{\parskip}{0.45em}
\setlength{\parindent}{0pt}
\onehalfspacing
\numberwithin{equation}{section}
\titleformat{\section}{\Large\bfseries\color{TitleBlue}}{\thesection}{0.6em}{}
\titleformat{\subsection}{\large\bfseries\color{AccentTeal}}{\thesubsection}{0.6em}{}
\newtcolorbox{overviewbox}{colback=SoftBlue,colframe=TitleBlue,boxrule=0.8pt,arc=2mm,left=3mm,right=3mm,top=2mm,bottom=2mm}
\newtcolorbox{notebox}{colback=SoftGray,colframe=AccentTeal,boxrule=0.6pt,arc=2mm,left=3mm,right=3mm,top=2mm,bottom=2mm}
\newcommand{\DocTitle}[3]{%
\begin{center}
{\LARGE \bfseries \color{TitleBlue} #1}\\[0.5em]
{\large \color{AccentTeal} #2}\\[0.8em]
{\normalsize #3}
\end{center}
\vspace{0.5em}}
\newcommand{\ProjectTOC}{\vspace{0.5em}{\hypersetup{linkcolor=TitleBlue}\tableofcontents}\vspace{0.8em}}
\allowdisplaybreaks

\newcolumntype{Y}{>{\RaggedRight\arraybackslash}X}

\begin{document}
\DocTitle{Project Plan}{PINN\_M3DC1}{Physics-informed surrogate architecture for M3D-C1-shaped nonlinear resistive-MHD workflows}

\begin{overviewbox}
\textbf{Purpose.} This is the living roadmap for \texttt{PINN\_M3DC1}. The project begins with a verified two-dimensional reduced-MHD benchmark. One surrogate is then trained first with simulation data alone and then with the reduced-MHD equations and known conditions added to its loss. Genuine M3D-C1 exports enter only through a later adapter and require a new validation claim.
\end{overviewbox}

\begin{notebox}
\textbf{Current status.} Notes for Modules~1, 2, and~3 are complete. Module~2 combines simulation, checking, saving, and data grouping. Module~3 combines the former data-only and physics-informed surrogate modules because the baseline is simply the same model with the physics-loss weights set to zero. Validation and M3D-C1 transfer are now Modules~4 and~5. No MHD solver, data pipeline, or trained neural network is yet implemented.
\end{notebox}

\ProjectTOC

\clearpage

\section{Objective and claim boundary}
The objective is to determine whether physics-informed learning improves a parameterized neural surrogate for nonlinear resistive-MHD evolution under controlled data budgets. Version 1 uses independently generated synthetic arrays. It supports the claim that a PINN can approximate a stated reduced model; it does not establish fidelity to M3D-C1. A validated M3D-C1 surrogate requires genuine code outputs, complete metadata, held-out M3D-C1 cases, and renewed verification.

\section{Initial physics benchmark}
The periodic two-field model evolves magnetic flux $\psi$ and stream function $U$. Define
\begin{equation}
 [f,g]=f_xg_y-f_yg_x,\qquad J=-\nabla^2\psi,\qquad \omega=\nabla^2U .
\end{equation}
The frozen version-1 residuals are
\begin{align}
 R_\psi &= \partial_t\psi+[U,\psi]-\eta\nabla^2\psi,\\
 R_\omega &= \partial_t\omega+[U,\omega]-[J,\psi]-\nu\nabla^2\omega .
\end{align}
The first nonlinear case is magnetic-island coalescence. It is small enough for independent implementation while retaining advection, current-sheet formation, resistive reconnection, dissipation, and parameter-dependent time evolution.

\section{Module roadmap}
\begin{singlespace}
\footnotesize
\renewcommand{\arraystretch}{1.12}
\begin{longtable}{p{0.07\linewidth} p{0.34\linewidth} p{0.50\linewidth}}
\toprule
\textbf{No.} & \textbf{Module} & \textbf{Purpose}\\
\midrule
1 & Reduced Resistive-MHD Physics and Conventions & Freeze variables, signs, normalization, boundary conditions, invariants, and diagnostic definitions.\\
2 & Generate and Save MHD Simulation Data & Run checked island-coalescence cases, save the numerical arrays and settings, and divide complete simulations into training, validation, and test groups.\\
3 & Physics-Informed Surrogate & Use one model for both settings: $\mathcal L_{\rm data}$ alone is the baseline; adding the two PDE residuals, the initial condition, and the gauge makes it physics-informed.\\
4 & Validation, Diagnostics, and Benchmarking & Compare field accuracy, physical histories, residuals, generalization, runtime, and individual loss-term tests.\\
5 & M3D-C1 Export Adapter and Transfer Learning & Map genuine exported fields and metadata into the project contract without silently changing physical meaning.\\
\bottomrule
\end{longtable}
\normalsize
\end{singlespace}

\section{Data and model contract}
A synthetic case stores $(x,y,t)$, $\psi$, $U$, $J$, $\omega$, $\eta$, $\nu$, $A_0$, the initial conditions, and the solver settings. The neural-network task is
\begin{equation}
 (\psi_0,U_0,\eta,\nu,t)\longmapsto(\widehat\psi(t),\widehat U(t)),
\end{equation}
with $\widehat J=-\nabla^2\widehat\psi$ and $\widehat\omega=\nabla^2\widehat U$. Plot images are not training data because they do not contain the complete numerical field values.

\section{Controlled comparison}
\begin{center}
\begin{tabularx}{\textwidth}{Yccc}
\toprule
\textbf{Model} & \textbf{Data loss} & \textbf{Physics losses} & \textbf{Role}\\
\midrule
Reduced-MHD solver & -- & Discretized equations & Synthetic reference\\
Module 3, data-only setting & yes & no & Fair neural baseline\\
Module 3, physics-informed setting & yes & PDE/IC/gauge & Physics-information test\\
\bottomrule
\end{tabularx}
\end{center}
The two neural models use identical case splits, observations, and comparable capacity. Entire $(\eta,\nu,A_0)$ combinations are excluded from training. Interpolation and extrapolation are reported separately.

\section{Verification and validation}
The reference solver must pass derivative, bracket, periodicity, divergence, convergence, and energy-behavior checks before producing training data. Surrogate evaluation includes relative $L_2$ and $L_\infty$ errors for all fields, peak-current and current-sheet errors, magnetic and kinetic energy histories, reconnection metrics, PDE residuals, and periodicity. Runtime reporting includes reference generation, training, inference, and break-even query count.

\clearpage
\section{Implementation sequence}
\begin{enumerate}[leftmargin=2em]
\item Freeze the equations and signs.
\item Implement the solver in Python and MATLAB.
\item Produce one checked island-coalescence simulation and save it.
\item Generate more cases and divide complete simulations into training, validation, and test groups.
\item Train the Module~3 model using $\mathcal L_{\rm data}$ alone.
\item Turn on the two PDE residual losses, the initial-condition loss, and the gauge loss.
\item Test optional constraints one at a time, then run generalization and cost analysis.
\item Add and validate an adapter only when genuine M3D-C1 arrays are available.
\end{enumerate}

\section{Acceptance milestones}
\textbf{M1 -- synthetic solver.} One command writes a reproducible HDF5 case and diagnostics; refinement and dissipation checks pass.

\textbf{M2 -- fair learning experiment.} The data-only and physics-informed settings of Module~3 share cases, samples, model assumptions, and scoring.

\textbf{M3 -- held-out validation.} Field and physics metrics are reported on unseen parameter combinations, with interpolation separated from extrapolation.

\textbf{M4 -- M3D-C1 transfer.} Exported variables, mesh, normalization, closures, and boundary conditions are mapped explicitly; the full test suite is rerun before an M3D-C1 surrogate claim is made.

\end{document}
